\documentclass[UTF8,a4paper,11pt]{ctexart}

\usepackage{amsmath,amssymb}
\usepackage{geometry}
\geometry{margin=2.5cm}
\usepackage{booktabs}
\usepackage{hyperref}
\usepackage{enumitem}

\title{控制系统：框图化简与 $G(s)$ 的含义}
\author{}
\date{}

\begin{document}
\maketitle

这张表总结了\textbf{控制系统框图（Block Diagram）等效变换的 10 条基本规则}，用于将复杂控制系统逐步化简成单个传递函数。

\bigskip\noindent\rule{\linewidth}{0.4pt}\bigskip

\subsection{1. 串联（Cascade）}

原框图：
\[
X \rightarrow G_1 \rightarrow G_2 \rightarrow Y
\]

等效：
\[
G_{\mathrm{eq}}=G_1G_2
\]

\textbf{规律：} 串联模块直接相乘。

\bigskip\noindent\rule{\linewidth}{0.4pt}\bigskip

\subsection{2. 并联（Parallel）}

原框图：
\[
Y=G_1X+G_2X
\]

等效：
\[
G_{\mathrm{eq}}=G_1+G_2
\]

\textbf{规律：} 并联模块直接相加（减号则相减）。

\bigskip\noindent\rule{\linewidth}{0.4pt}\bigskip

\subsection{3. 提取公共模块}

将
\[
Y=G_1X+G_2X
\]
改写为
\[
Y=G_2\left(X+\frac{G_1}{G_2}X\right)
\]

\textbf{作用：} 为进一步化简创造条件。

\bigskip\noindent\rule{\linewidth}{0.4pt}\bigskip

\subsection{4. 消除反馈环（最重要）}

负反馈系统：
\[
Y=G_1(X-G_2Y)
\]

化简得到：
\[
G_{\mathrm{eq}}=\frac{G_1}{1+G_1G_2}
\]

正反馈则为：
\[
G_{\mathrm{eq}}=\frac{G_1}{1-G_1G_2}
\]

这是控制系统中最常用的公式。

\bigskip\noindent\rule{\linewidth}{0.4pt}\bigskip

\subsection{5. 从反馈回路中移出模块}

用于把反馈通道中的模块移到前向通道或反之。

核心思想：
\[
G_1(X\pm G_2Y)
\]

通过调整结构保持输入输出关系不变。

\bigskip\noindent\rule{\linewidth}{0.4pt}\bigskip

\subsection{6. 调整加法点顺序}

原式：
\[
Z=W\pm X\pm Y
\]

无论先加哪个信号：
\[
(W\pm X)\pm Y=(W\pm Y)\pm X
\]

结果相同。

\bigskip\noindent\rule{\linewidth}{0.4pt}\bigskip

\subsection{7. 加法点前移（跨过模块）}

原结构：
\[
Z=GX\pm Y
\]

如果把求和点移到模块前面，$Y$ 必须经过 $\frac{1}{G}$ 补偿。

原因：模块 $G$ 会同时作用于两个输入。

\bigskip\noindent\rule{\linewidth}{0.4pt}\bigskip

\subsection{8. 加法点后移（跨过模块）}

原结构：
\[
(X\pm Y)\rightarrow G
\]

变换后：
\[
GX \pm GY
\]

即：
\[
G(X\pm Y)=GX\pm GY
\]

\bigskip\noindent\rule{\linewidth}{0.4pt}\bigskip

\subsection{9. 分支点前移}

若分支点从模块后移到模块前，需要在新分支上增加一个 $G$。

原因：原分支信号是 $GX$。

\bigskip\noindent\rule{\linewidth}{0.4pt}\bigskip

\subsection{10. 分支点后移}

若分支点从模块前移到模块后，需要增加 $\frac{1}{G}$ 补偿。

\bigskip\noindent\rule{\linewidth}{0.4pt}\bigskip

\section{考试最常考的三条}

\begin{enumerate}
  \item \textbf{串联相乘}
  \[
  G=G_1G_2
  \]

  \item \textbf{并联相加}
  \[
  G=G_1+G_2
  \]

  \item \textbf{反馈化简}
  \[
  G=\frac{G_1}{1+G_1H} \quad \text{（负反馈）}
  \]
  \[
  G=\frac{G_1}{1-G_1H} \quad \text{（正反馈）}
  \]
\end{enumerate}

如果你正在学习\textbf{控制系统、自动控制原理或信号与系统}，掌握这 10 条规则后，绝大多数框图化简题都可以通过“串联 $\rightarrow$ 并联 $\rightarrow$ 消反馈 $\rightarrow$ 调整加法点和分支点”的顺序一步步求出总传递函数。

\bigskip\noindent\rule{\linewidth}{0.4pt}\bigskip

这是很多人学控制系统时最困惑的地方。

因为教材一上来就写：
\[
G(s)=\frac{Y(s)}{X(s)}
\]

然后开始推公式，但从来不解释：

\begin{quote}
\textbf{这个 $G(s)$ 到底是什么东西？}
\end{quote}

其实最直观的理解是：

\begin{quote}
\textbf{$G(s)$ 描述了一个系统对输入信号的反应规律。}
\end{quote}

你可以暂时忘掉拉普拉斯变换，把 $G(s)$ 看成一个“输入$\rightarrow$输出转换器”。

\bigskip\noindent\rule{\linewidth}{0.4pt}\bigskip

\subsection{先从最简单的例子开始}

假设有一个黑盒：

\begin{verbatim}
输入 x
 ↓
[黑盒]
 ↓
输出 y
\end{verbatim}

如果：

\begin{verbatim}
输入 2    输出 10
输入 3    输出 15
输入 4    输出 20
\end{verbatim}

你会发现：
\[
y=5x
\]

那么这个黑盒的增益就是
\[
G=5
\]

意思是：

\begin{quote}
输入放大 5 倍输出。
\end{quote}

\bigskip\noindent\rule{\linewidth}{0.4pt}\bigskip

\subsection{电机例子}

假设：

\begin{verbatim}
输入电压
 ↓
电机
 ↓
转速
\end{verbatim}

实验发现：

\begin{center}
\begin{tabular}{cc}
\toprule
电压 (V) & 转速 (rpm) \\
\midrule
1 & 100 \\
2 & 200 \\
3 & 300 \\
\bottomrule
\end{tabular}
\end{center}

那么：
\[
\text{转速}=100\times \text{电压}
\]

于是可以写：
\[
G=100
\]

$G$ 表示：

\begin{quote}
电机把电压转换成转速的能力。
\end{quote}

\bigskip\noindent\rule{\linewidth}{0.4pt}\bigskip

\subsection{为什么后来变成 $G(s)$？}

现实系统不是立刻响应。

例如：你给电机 10V。

理想情况：

\begin{verbatim}
t=0    转速=1000rpm
\end{verbatim}

现实情况：

\begin{verbatim}
t=0      0rpm
t=1s   500rpm
t=2s   800rpm
t=3s   950rpm
t=4s 1000rpm
\end{verbatim}

速度是慢慢爬上去的。

所以系统不仅有：

\begin{itemize}
  \item 放大倍数
\end{itemize}

还有：

\begin{itemize}
  \item 响应速度
  \item 延迟
  \item 振荡
\end{itemize}

这些特性。

一个数字 $G$ 已经描述不了。

于是需要 $G(s)$。

\bigskip\noindent\rule{\linewidth}{0.4pt}\bigskip

\subsection{一阶系统}

最常见的电机模型：
\[
G(s)=\frac{K}{Ts+1}
\]

其中：

\begin{itemize}
  \item $K$ = 最终放大倍数
  \item $T$ = 响应快慢
\end{itemize}

例如：
\[
G(s)=\frac{10}{2s+1}
\]

意思：

\begin{itemize}
  \item 最终输出是输入的 10 倍
  \item 需要大约 2 秒才能跟上
\end{itemize}

所以 $G(s)$ 不仅告诉你放大多少，还告诉你：

\begin{quote}
\textbf{多久达到这个结果。}
\end{quote}

\bigskip\noindent\rule{\linewidth}{0.4pt}\bigskip

\subsection{汽车的例子}

油门：

\begin{verbatim}
输入：油门开度
\end{verbatim}

车速：

\begin{verbatim}
输出：车速
\end{verbatim}

现实情况，踩到底后：

\begin{verbatim}
0s    0km/h
2s   20km/h
5s   60km/h
10s 100km/h
\end{verbatim}

汽车不会瞬间达到 100km/h。

因此汽车可以抽象成：

\begin{verbatim}
油门
 ↓
G(s)
 ↓
车速
\end{verbatim}

这里的 $G(s)$ 就是：

\begin{quote}
这辆车的动力学特性。
\end{quote}

\bigskip\noindent\rule{\linewidth}{0.4pt}\bigskip

\subsection{为什么控制系统全是 $G_1$、$G_2$、$G_3$？}

因为工程师喜欢把复杂系统拆开。

例如无人机：

\begin{verbatim}
控制器
 ↓
电机
 ↓
机体
 ↓
传感器
\end{verbatim}

对应：

\begin{verbatim}
Gc(s)
 ↓
Gm(s)
 ↓
Gp(s)
 ↓
H(s)
\end{verbatim}

其中：

\begin{itemize}
  \item $G_c$ 控制器
  \item $G_m$ 电机
  \item $G_p$ 机体
  \item $H$ 传感器
\end{itemize}

每个模块都表示：

\begin{quote}
输入经过这个部分后会怎样变化。
\end{quote}

\bigskip\noindent\rule{\linewidth}{0.4pt}\bigskip

\subsection{最直观的一句话}

学框图时，你可以把 $G(s)$ 直接理解成：

\begin{quote}
\textbf{某个设备、某个部件、某个过程的“行为说明书”。}
\end{quote}

例如：

\begin{center}
\begin{tabular}{ll}
\toprule
$G(s)$ & 代表 \\
\midrule
电机 $G(s)$ & 电压$\rightarrow$转速 \\
汽车 $G(s)$ & 油门$\rightarrow$车速 \\
空调 $G(s)$ & 功率$\rightarrow$温度 \\
飞机 $G(s)$ & 舵角$\rightarrow$姿态 \\
音响 $G(s)$ & 输入电压$\rightarrow$声音 \\
\bottomrule
\end{tabular}
\end{center}

框图里的每个方块，其实都在回答：

\begin{quote}
给它一个输入，它会怎样随时间产生输出？
\end{quote}

而整个框图化简，就是把许多个“小行为说明书”合并成一个“大行为说明书”，最终得到：
\[
\frac{Y(s)}{R(s)}
\]

也就是：

\begin{quote}
从输入到最终输出，整个系统到底会怎么表现。
\end{quote}

\end{document}
